
\chapter{Реестр обязательств и~события}\label{ch:internal-ledger}

Шлюз отвечает продавцу за~сотни миллисекунд. К~этому моменту ещё не~ясно:
произойдёт ли подтверждение (\textit{capture}~--- вторая стадия двухстадийной
  авторизации, на~которой ранее зарезервированная сумма списывается
фактически)~\cite{stripe-manual-capture,adyen-glossary}, когда придёт расчёт
(см.~раздел~\ref{sec:tx-clearing}), какой будет комиссия эквайера
и~не~превратится ли успешная продажа через месяц в~чарджбэк. Между шлюзом
и~сверкой нужен ещё один компонент~--- внутренний реестр обязательств,
в~котором продукт фиксирует собственные денежные обязательства и~привязывает их
к~событиям платёжного
потока~\cite{visa-core-rules,mc-tpr,stripe-manual-capture}.

\section{Шлюз не~заменяет учёт}\label{sec:ledger-why}

Шлюз отвечает на~вопрос \enquote{что произошло на~входе в~систему}: пришёл
запрос, сработала проверка риска, эквайер вернул одобрение или отказ, продавцу
ушёл вебхук. Для интеграции этого достаточно, для финансового учёта~--- нет.

Во-первых, статус шлюза не~равен фактическому положению денег. Авторизация
со~статусом \enquote{approved} означает лишь, что эмитент зарезервировал сумму
или одобрил её к~дальнейшей обработке. Подтверждение, клиринг
(\textit{clearing}~--- обмен данными о~транзакциях между банками-участниками),
расчёт (\textit{settlement}~--- фактическое перемещение денежных средств между
ними; см.~раздел~\ref{sec:tx-clearing}), отмены и~корректировки в~правилах
карточных схем идут отдельными стадиями~\cite{mc-switching-services}. Пока
сумма не~перешла в~доступный остаток, считать её готовой к~выплате
(\textit{payout}~--- перевод доступного остатка с~внутреннего баланса продавца
на~банковский счёт~\cite{stripe-payouts})
рано~\cite{visa-core-rules,mc-tpr,stripe-manual-capture,stripe-balances}.

Во-вторых, за~одним внешним статусом скрывается несколько денежных последствий.
Успешная транзакция в~маркетплейсе порождает отдельные обязательства перед ТСП,
платформой и~эквайером. Хранить это одной строкой \texttt{payment.status =
succeeded}~--- значит рано или поздно потерять связь между движением денег
и~жизненным циклом операции.

В~маркетплейсе разрыв виден ещё чётче: покупатель платит одну сумму, а~из~неё
возникают выплата продавцу, комиссия платформы и~удержание на~доставку, если
она есть. Если платформа гарантирует возврат, добавляется резерв под возвраты.
В~шлюзе всё это одна строка \texttt{payment = 5000}, тогда как в~реестре~---
проводки с~разными сроками признания и~разными правилами отмены.

\begin{figure}[tbp]
  \centering
  \includefiguresvg[width=.8\linewidth]{assets/figures/ch34-ledger/gateway-vs-ledger}
  \caption{Шлюз против реестра обязательств: один платёж, четыре проводки.}\label{fig:gateway-vs-ledger}
\end{figure}

Шлюз из~раздела~\ref{sec:gw-anatomy} отвечает за~приём и~маршрутизацию.
\highlight{Реестр отвечает за~состояние денег внутри продукта: что уже
заработано, что зарезервировано, что можно выплатить продавцу}.

\section{Двойная запись: минимум для платёжного продукта}\label{sec:ledger-double-entry}

\textbf{Счёт}~--- именованный накопитель денежных значений:
\texttt{cardholder\_receivable}, \texttt{merchant\_pending},
\texttt{platform\_fee\_earned}. \textbf{Проводка} (бухгалтерская запись)~---
пара \enquote{дебет счёта~A, кредит счёта~B на~одну и~ту же сумму}. Один счёт
увеличивается, другой уменьшается, в~зависимости от~типа. \textbf{Двойная
запись} (double-entry bookkeeping) фиксирует каждое денежное событие минимум
одной парной проводкой, при~которой сумма дебетов равна сумме кредитов.
Расхождение сумм означает ошибку в~журнале: инвариант равенства встроен в~саму
схему учёта.

Двойная запись отличает реестр обязательств от~обычного журнала. Журнал событий
фиксирует \enquote{покупатель оплатил заказ}; реестр выводит из~этого парные
проводки, у~которых сумма дебетов всегда равна сумме кредитов.

Аудиторский след при~двойной записи сам себя проверяет: \enquote{потерять}
деньги без видимого расхождения нельзя. Обязательства и~фактические остатки
учитываются раздельно (pending против available,
см.~раздел~\ref{sec:ledger-balances}). Банковская выписка устроена по~той же
двойной записи со~стороны банка, и~сопоставление двух согласованных
представлений надёжнее сравнения журнала с~одной агрегированной цифрой.

\section{Продуктовый реестр обязательств и~бухгалтерский учёт банка}\label{sec:ledger-vs-accounting}

Внутренний реестр платёжного продукта и~бухгалтерский учёт кредитной
организации, внутри которой он работает,~--- две разные системы. Их легко
перепутать, особенно если продукт встроен в~банк.

\textbf{Продукт}~--- платёжный сервис, который видит покупатель и~продавец:
эквайринговый продукт банка, PSP, агрегатор, маркетплейс, кошелёк. Это уровень
бизнес-логики (авторизация, capture, выплата ТСП, чарджбэк, комиссия
платформы). \textbf{Банк}~--- кредитная организация, на~счетах которой реально
движутся деньги по~809-П: у~банковского продукта это сам банк-оператор, у~PSP
или платформы~--- банк-эквайер или банк-партнёр. Эти два учётных уровня могут
совпадать (продукт живёт внутри банка) или быть разделены (продукт работает
поверх партнёрского банка). Реестр обязательств~--- инструмент продукта; счета
809-П~--- банка.

\textbf{Продуктовый реестр обязательств}~--- внутренний инструмент платёжного
продукта. Он живёт в~базах данных продукта, использует произвольную нотацию
счетов (\texttt{merchant\_pending}, \texttt{platform\_fee\_earned},
\texttt{refund\_payable}), считает балансы партнёров~--- продавцов,
плательщиков, платформ~--- и~отвечает на~вопросы продукта: сколько мы должны
ТСП, сколько уже признано комиссией, какой резерв удержан под возвраты. Наружу
как бухгалтерский отчёт он не~публикуется.

\textbf{Бухгалтерский учёт банка по~РСБУ} ведётся кредитной организацией
по~стандартам РСБУ на~основе Плана счетов кредитных организаций, утверждённого
Положением Банка России от~24.11.2022~№~809-П \enquote{О~Плане счетов
  бухгалтерского учёта для кредитных
организаций}~\cite{cbr-809p}.\footnote{В~силе с~01.01.2023; заменило 579-П.}
План счетов задаёт фиксированную нумерацию, где \enquote{30102} обозначает
корреспондентские счета кредитных организаций в~Банке России,
\enquote{40817}~--- счета физических лиц, \enquote{40702}~--- счета
коммерческих организаций, а~\enquote{47422}~--- \enquote{Обязательства
по~прочим операциям}, на~котором, в~частности, учитываются суммы клиентских
платежей в~пути. Банк обязан отражать платёжные операции на~этих счетах
в~установленном порядке: содержимое подаётся в~формальной отчётности Банку
России и~проверяется аудитом.

Продуктовый реестр даёт детальный взгляд на~одну операцию по~её жизненному
циклу внутри продукта. Бухгалтерский учёт показывает общий срез по~банковским
счетам и~обязательным группировкам. Каждая запись продуктового реестра
в~конечном счёте порождает одну или несколько бухгалтерских проводок на~счетах
809-П, а~обратное верно не~всегда, потому что бухгалтерия агрегирует
продуктовые записи и~хранит их в~свёрнутом виде. Для платёжного продукта
в~банке соответствие между внутренними счетами продуктового реестра и~счетами
809-П фиксируется явно: сверка тогда идёт по~этому соответствию, без
дублирования двойной записи в~обеих системах и~без пересборки бухгалтерии
заново.

Если платёжный продукт работает вне банка (например, в~составе небанковской
платформы), связь смещается: бухгалтерия по~РСБУ ведётся у~банка-партнёра или
нескольких партнёров; продуктовый реестр связан с~ней через интерфейсные
документы, реестры и~банковские выписки. В~обоих случаях соответствие между
внутренним реестром и~Планом счетов нужно закладывать в~архитектуру с~самого
начала~--- иначе получаются два независимых параллельных учёта, которые потом
не~сойдутся.

\section{Три класса реализации реестра}\label{sec:ledger-implementation-classes}

Реестр обязательств можно построить тремя способами: двойную запись
из~раздела~\ref{sec:ledger-double-entry} гарантирует прикладной код, сторонний
сервис или сама СУБД. Выбор задаёт стоимость поддержки, потолок
производительности и~цену миграции на~следующую систему.

\subsection*{Реестр как код поверх универсальной СУБД}

Проводки хранятся в~обычной реляционной или документной базе, а~инвариант
двойной записи обеспечивает прикладной код. Именно эту схему реализует листинг
\texttt{posting.py} из~раздела~\ref{sec:ledger-events}. База остаётся
универсальной (PostgreSQL, MySQL, ClickHouse, MongoDB), учётная логика живёт
в~сервисах продукта. Гибкость максимальна: можно свободно менять схему счетов,
правила начисления комиссии, переписывать конечный автомат переходов. Цена~---
дисциплина: ошибка в~одной из~ветвей кода способна оставить незакрытое
расхождение, а~автотесты ловят только явно описанные сценарии.

\subsection*{Реестр как сторонний сервис}

Внешний API, который продукт вызывает для записи проводок и~запроса балансов.
Инварианты двойной записи держит сервис, продукт отправляет ему события,
не~зная, как внутри устроено хранилище. Modern Treasury
Ledgers~\cite{modern-treasury-ledgers}~--- облачный реестр на~правах сервиса
(\textit{managed}) с~декларируемой неизменяемостью (\textit{immutability}):
запись, опубликованная в~журнал, не~редактируется и~не~удаляется, а~исправления
оформляются встречными проводками.\footnote{В~октябре 2021~года Modern Treasury
  закрыл первичный раунд Series C на~\$85\,млн при~валуации более~\$2\,млрд,
  а~в~марте 2022~года~--- расширение того же раунда на~\$50\,млн от~SVB Capital
  и~Salesforce Ventures; суммарный Series C составил
\$135\,млн~\cite{modern-treasury-funding}.} Formance
Ledger~\cite{formance-ledger}~--- открытый исходный код (Apache~2.0)
с~собственным DSL~\textit{numscript}, на~котором проводки описываются как
программа, а~не как набор инструкций ORM.\footnote{Formance прошёл батч
  Summer~2021 Y~Combinator с~\textit{pre-seed}-раундом на~\$3{,}1\,млн
(Y~Combinator, Hoxton Ventures, Frst)~\cite{formance-yc}.} Blnk
Finance~\cite{blnk-finance}~--- модульный реестр с~открытым исходным кодом
в~той же нише: двойная запись, неизменяемый журнал, встроенные сверка
и~проводки в~полёте (\textit{inflight}). Подходит, когда продукт не~хочет
владеть инфраструктурой реестра и~готов платить за~SLA или брать на~себя
эксплуатацию открытого сервиса у~себя.

\subsection*{Реестр как сама база данных}

TigerBeetle~\cite{tigerbeetle-debit-credit}~--- OLTP-СУБД с~зашитой в~схему
данных моделью двойной записи: каждый \texttt{transfer} обязан содержать
\texttt{debit\_account\_id} и~\texttt{credit\_account\_id} с~равной суммой,
любая попытка нарушить инвариант отклоняется на~уровне БД, а~не на~уровне
приложения. Журнал работает только на~дозапись (\textit{append-only}): проводку
нельзя удалить или переписать, обратные движения оформляются отдельными
записями. Декларируемая пропускная способность~--- до~1\,000\,000 переводов
в~секунду на~одиночном узле. За~эту скорость и~встроенную целостность платят
узкой специализацией: TigerBeetle не~годится для произвольных запросов,
бизнес-логики поверх проводок или сложной аналитики. В~архитектуре приложения
он оказывается рядом с~основной СУБД, а~не вместо неё.

\highlight{Выбор класса реестра~--- архитектурное решение того же порядка, что
  выбор основной СУБД: переход с~одного класса на~другой задним числом обходится
в~отдельный проект миграции данных и~учётной политики}.

\section{События: авторизация, подтверждение, возврат, спор}\label{sec:ledger-events}

В~карточном потоке реестр обрабатывает события авторизации, подтверждения,
возврата, чарджбэка и~начисления комиссии. Каждое порождает собственный набор
проводок и~не~сводится к~общему статусу \enquote{платёж
успешен}~\cite{mc-tpr,mc-chargeback-guide,adyen-glossary}. Реестр атомарно
записывает событие в~журнал и~сразу порождает проводки, поэтому продуктовая
система видит две взаимосвязанные картины одной операции: \textit{что
произошло} и~\textit{как это отразилось в~деньгах}.

Авторизация на~5\,000~руб. ещё не~создаёт выручку ТСП, но фиксирует ожидаемое
обязательство. Подтверждение переводит обязательство в~признанный доход;
до~расчёта эти деньги ещё не~на~балансе, доступном к~выплате. Возврат запускает
обратное движение уже после подтверждения; платёж, не~дошедший
до~подтверждения, обычно отменяют (возврат для него
не~применяется)~\cite{adyen-glossary,stripe-manual-capture}. Чарджбэк~---
самостоятельное событие, инициируемое банком-эмитентом, которое живёт
по~правилам платёжной системы и~несёт собственные сроки и~комиссии
(см.~раздел~\ref{sec:disputes-lifecycle})~\cite{adyen-glossary,mc-chargeback-guide}.

\practicenote{%
  Кодируйте возврат и~чарджбэк как отдельные типы операций.
  Для клиента результат похож, но для продукта это разные
  процессы: возврат инициирует продавец, а~чарджбэк приходит извне
  и~живёт в~цикле оспаривания.
}

\highlight{Событие отвечает за~причинность, проводка~--- за~денежный смысл.
  Попытка восстановить одно из~другого задним числом лишает продукт
объяснимости}.\footnote{Тот же подход~--- системы-источники как конечные
  автоматы, чьи переходы разворачиваются реестром в~проводки~--- описан
  на~архитектурном уровне в~инженерном обзоре собственного реестра
  Stripe~\cite{stripe-ledger-blog}: декларируется около~5\,млрд событий реестра
  в~сутки и~объяснимость (\textit{explainability}) движения денег выше
99{,}9999\,\%.}

Переходы образуют конечный автомат: авторизация переходит в~подтверждение или
в~отмену (\textit{void}~--- аннулирование авторизации до~её подтверждения;
  финансовой проводки в~сети
не~порождает~\cite{adyen-glossary,stripe-manual-capture}), но не~наоборот;
подтверждение переходит в~\enquote{рассчитанное} состояние, после чего возможны
возврат или чарджбэк; отмена допустима только до~подтверждения. Возврат после
расчёта порождает обратное движение средств, причём частичный возврат уменьшает
обязательство, не~отменяя его. Чарджбэк инициируется извне и~возможен после
расчёта; он открывает отдельную цепочку~--- повторное представление
(\textit{representment}~--- ответ ТСП с~доказательствами легитимности операции;
у~Mastercard этот шаг называется \textit{second presentment}), доарбитраж
(\textit{pre-arbitration}~--- спор после повторного представления) и~арбитраж
(\textit{arbitration}~--- финальное разрешение спора платёжной системой);
подробнее~---
раздел~\ref{sec:disputes-lifecycle}~\cite{visa-core-rules,mc-chargeback-guide}.

Каждый переход порождает проводку. Сумма обязательства перед ТСП движется
по~трёхфазной модели
\texttt{pending}~$\to$~\texttt{earned}~$\to$~\texttt{available}, где
\texttt{pending}~--- зафиксированное, но ещё не~признанное обязательство (после
авторизации до~подтверждения); \texttt{earned}~--- признанная выручка ТСП
(после подтверждения, ещё до~фактического расчёта); \texttt{available}~---
доступный к~выплате остаток (после
расчёта)~\cite{stripe-balances,adyen-balances}. Авторизация создаёт
\texttt{pending}, подтверждение переводит \texttt{pending} в~\texttt{earned},
отмена обнуляет \texttt{pending}, расчёт переводит \texttt{earned}
в~\texttt{available}. Возврат дебетует \texttt{available} и~кредитует
\texttt{refund\_payable}, а~чарджбэк дебетует \texttt{available} или
\texttt{reserve} и~открывает отдельный учётный цикл спора.

\importantnote{%
  Состояние \texttt{pending} в~реестре и~\emph{холд на~счёте
  держателя}, который ставит эмитент после авторизации
  (см.~раздел~\ref{sec:tx-authorization}),~--- два разных таймера.
  Реестровое \texttt{pending} управляется внутренней логикой
  продукта~--- моментами подтверждения, расчёта, отмены; холд
  эмитента на~авторизации (\textit{auth-hold}) живёт по~срокам
  Visa/MC и~снимается автоматически по~истечении срока,
  установленного правилами платёжной системы для категории
  операции, даже если в~реестре проводка остаётся.
  Рассинхронизацию между ними обрабатывает сценарий
  \emph{orphan auth} (разбирается ниже
  в~разделе про краевые случаи).
}

Без явно описанного автомата команда кодирует переходы в~виде разрозненных
условий, и~через год никто не~помнит, допустим ли возврат после частичного
чарджбэка.

Продавец маркетплейса продаёт товар за~5\,000~руб. при~комиссии платформы
10\,\%; на~этом примере проводки выглядят так.

\textbf{Auth (день~0).} Когда эмитент одобряет авторизацию, реестр создаёт две
проводки: дебет \texttt{cardholder\_receivable} 5\,000~руб. и~кредит
\texttt{merchant\_pending} 5\,000~руб. Деньги ещё не~у~платформы, но
обязательство зафиксировано.

\textbf{Подтверждение (день~0, через 30~мин).} Когда продавец подтверждает
отгрузку, реестр переводит pending в~earned проводками: дебет
\texttt{merchant\_pending} 5\,000~руб., кредит \texttt{merchant\_earned}
4\,500~руб. и~кредит \texttt{platform\_fee\_earned} 500~руб. С~этого момента
платформа признала свою комиссию.

\textbf{Расчёт (день~2).} Когда эквайер перечисляет средства, реестр переводит
earned в~available проводками: дебет \texttt{merchant\_earned} 4\,500~руб.
и~кредит \texttt{merchant\_available} 4\,500~руб. После этого продавец вправе
запросить выплату.

\textbf{Частичный возврат (день~5).} Когда покупатель возвращает часть товара
на~2\,000~руб., реестр оформляет это как дебет \texttt{merchant\_available}
2\,000~руб. и~кредит \texttt{refund\_payable} 2\,000~руб. Комиссия платформы
не~пересчитывается (политика маркетплейса); если бы пересчитывалась, отдельная
проводка уменьшила бы \texttt{platform\_fee\_earned} на~200~руб.

Каждый переход оставляет аудиторский след: сумма дебетов равна сумме кредитов,
а~баланс любого счёта восстанавливается обходом журнала проводок.

Знак сальдо следует РСБУ-классификации счёта. Активный счёт
(\texttt{cardholder\_receivable}~--- требование к~плательщику, \enquote{где
деньги}) накапливает дебет и~даёт отрицательный остаток. Пассивные счета
(обязательства перед продавцом и~плательщиком, признанная выручка~---
\enquote{откуда деньги}) накапливают кредит и~дают положительный.

Проводка~--- неизменяемая запись \texttt{(дебет, кредит, сумма)}; реестр~---
журнал таких проводок и~операция запроса сальдо; событие платёжного потока~---
функция, формирующая группу проводок.

Минимальный API реестра~--- три типа и~четыре функции событий. \texttt{Leg}
хранит одну ногу проводки и~в~конструкторе отвергает нулевые и~отрицательные
суммы: пустая проводка чаще означает ошибку в~коде, чем правомерный сценарий,
и~запрет на~конструкторе дешевле, чем восстановление инварианта по~журналу
постфактум. \texttt{Ledger.post} принимает переменное число ног и~записывает их
одной операцией: событие \texttt{confirm} порождает две проводки сразу,
и~журнал не~остаётся в~полусобранном состоянии при~сбое между ними.
\texttt{balance} восстанавливает сальдо обходом журнала вместо хранения
текущего остатка рядом с~проводками: медленнее в~эксплуатации, зато инвариант
\enquote{сумма дебетов равна сумме кредитов} проверяется тем же кодом, который
выводит сальдо, без отдельной фоновой ребалансировки. В~рабочей нагрузке прямой
обход журнала заменяется производными структурами: snapshot-таблицей текущего
остатка по~счёту, материализованным представлением \texttt{balance}
с~инкрементальным обновлением или триггером накопления при~записи проводки. Сам
журнал остаётся источником истины, производные структуры~--- кеш с~явной
процедурой пересборки из~журнала.

\codefile{Python}{samples/ch34-ledger/posting.py}

Внутренние счета реестра не~совпадают один к~одному со~счетами 809-П на~стороне
банка-эквайера (или банка-партнёра, если речь о~небанковском продукте).

\texttt{merchant\_pending} соответствует \texttt{None}: авторизация~---
коммуникация со~схемой, а~не движение по~счетам банка, и~в~809-П
авторизационный hold балансовой проводкой не~оформляется. Продуктовый реестр
это состояние различает, банковский учёт~--- ещё нет. Соответственно и~дебет
\texttt{cardholder\_receivable} в~реестре не~синхронен с~hold'ом, который
держит эмитент на~счёте плательщика: hold снимается по~истечении срока,
установленного правилами сети для категории
операции~(см.~раздел~\ref{sec:tx-authorization}), проводка реестра~--- только
событием клиринга или отмены.

После клиринга появляется реальная пара 30233\,/\,47422 (30233~---
  \enquote{Незавершённые расчёты с~операторами услуг платёжной инфраструктуры
  и~операторами по~переводу денежных средств}, активный счёт; 47422~---
\enquote{Обязательства по~прочим операциям}, пассивный)~\cite{cbr-809p}: банк
фиксирует требование к~оператору платёжной инфраструктуры и~встречное
обязательство перед ТСП. \texttt{merchant\_earned}
и~\texttt{merchant\_available} отображаются в~один и~тот же 47422~---
аналитические срезы одного и~того же обязательства до~выплаты ТСП.
\texttt{refund\_payable} живёт там же, но с~противоположным знаком движения.
\texttt{platform\_fee\_earned} уходит на~70601 (\enquote{Доходы}~---
операционные доходы кредитной организации~\cite{cbr-809p}), если речь
о~банке-партнёре или эквайере; в~небанковском продукте этой проводки на~809-П
не~возникает вовсе, доход живёт в~учёте платформы. Конкретное соответствие
утверждается учётной политикой кредитной организации и~зависит от~схемы
расчётов.

\begin{figure}[tbp]
  \centering
  \includefiguresvg[width=\linewidth]{assets/figures/ch34-ledger/balance-event-grid}
  \caption{Балансы по~горизонтали, события по~вертикали.}%
  \label{fig:balance-event-grid}
\end{figure}

\section{Балансы, комиссии и~правила перехода}\label{sec:ledger-balances}

\textbf{Pending} и~\textbf{available}~--- состояния одного и~того же основного
баланса. Pending~--- сумма, недоступная к~выплате: она зафиксирована между
авторизацией и~подтверждением, ожидает клиринга или удерживается до~даты
ближайшей выплаты. Available~--- та часть, которой платформа готова
распоряжаться: перечислить продавцу, зачесть в~счёт комиссии или использовать
для внутреннего неттинга (\textit{netting}~--- взаимозачёт встречных требований
до~фактического перевода средств). Деньги переходят из~pending в~available
по~мере прохождения этапов жизненного цикла
операции~\cite{stripe-balances,adyen-balances}.

\textbf{Резерв (\textit{reserve})}~--- отдельный балансовый механизм, а~не ещё
одно состояние того же баланса. Резерв хранит сумму, изъятую из~доступного
остатка по~правилу риска, и~подчиняется собственным правилам пополнения,
разморозки и~закрытия. Типичные разновидности:
\begin{itemize}
  \item \textit{refund reserve}~--- удержание под будущие возвраты,
    рассчитанное от~объёма возвратной активности ТСП;
  \item минимальный остаток (\textit{minimum balance})~--- неснижаемая сумма
    на~балансе ТСП для покрытия будущих споров и~сборов;
  \item фиксированный резерв (\textit{holdback}~--- удержание фиксированной
    доли каждой транзакции до~окончания согласованного срока) под ожидаемые
    списания~\cite{stripe-balances,stripe-minimum-balances,adyen-balances}.
\end{itemize}

Без явного моделирования переходов pending~$\to$~earned~$\to$~available продукт
перестаёт отличать \enquote{деньги есть} от \enquote{деньги когда-нибудь
придут}. Отсюда и~расхождения между панелью продавца и~банковской выпиской.

Комиссия платформы редко признаётся одновременно с~авторизацией: её признают
при~подтверждении или расчёте, когда ясно, что операция состоялась. Эквайерские
сборы, сетевые комиссии, валютная наценка и~комиссии за~чарджбэк появляются
по~собственным графикам, отличным от~графика основной суммы. Реестр ведёт
начисления в~два прохода: сначала признаёт ожидаемую комиссию, затем
корректирует её по~фактическому расчётному
файлу~\cite{visa-core-rules,mc-tpr,stripe-balances}.

То же касается резервов: резерв возвратов, minimum balance или любой другой
удерживаемый резерв описывается как конкретное правило перевода части available
в~отдельный баланс с~собственным сроком разморозки.

Основной поток балансов и~цикл оспаривания после чарджбэка показаны
на~рис.~\ref{fig:balance-event-grid}.

График выплат тоже моделируется в~реестре, а~не в~коде экспортёра банковских
файлов. Ежедневные, еженедельные и~ежемесячные циклы перечислений, ручные
перечисления и~задержки вида T+X~--- всё это правила перевода между внутренними
балансами~\cite{stripe-payouts}. Правило, не~выражённое в~реестре, повторяется
вручную в~процессе выплат и~бэк-офиса, и~единый источник истины
теряется~\cite{stripe-minimum-balances,adyen-balances}.

\section{Идемпотентность и~передача в~сверку}\label{sec:ledger-idempotency}

Идемпотентность на~API шлюза разобрана в~разделе~\ref{sec:gw-reliability}. Шлюз
присваивает операции стабильный внутренний идентификатор и~связывает с~ним все
внешние сообщения через таблицу соответствия с~STAN, RRN и~другими
идентификаторами процессинга~\cite{fis-iso8583-guide-2023}. Реестр принимает
только нормализованные внутренние события с~уникальным \path{event_id}.
Повторная доставка такого события возвращает прежний результат, не~порождая
новых проводок~\cite{stripe-webhooks}.

Запись события в~журнал и~применение к~балансам разделены на~два шага: сначала
фиксируется факт получения уведомления, затем событие атомарно меняет балансы.
После сбоя между шагами повторный проход доводит проводки до~конца, не~дублируя
их.

Сверка из~главы~\ref{ch:reconciliation} сопоставляет внутреннее состояние
продукта с~внешними файлами и~банковскими движениями. Для этого продукт обязан
сначала зафиксировать собственное состояние. Шлюз записывает протокольный факт
авторизации, реестр переводит этот факт в~денежные движения, а~сверка
сопоставляет результат с~TC33, IPM и~банковской
выпиской~\cite{visa-core-rules,mc-tpr,mc-switching-services}.

Без такого реестра сверка напрямую сопоставляет внешние файлы со~статусами
шлюза, и~картина бэк-офиса становится мутной: \enquote{вроде операция одобрена,
  но непонятно, это уже доход, ожидаемая выплата или просто авторизационная
блокировка}.

\needspace{4\baselineskip}
\subsection*{Краевые случаи: orphan auth и~гонка возврата с~чарджбэком}

Реестры, спроектированные под счастливый путь (\textit{happy path}), ломаются
на~незавершённой авторизации и~на~одновременном возврате со~спором. Оба
сценария требуют явной обработки в~коде реестра, иначе расхождения копятся
незаметно.

\textbf{Orphan auth.} Когда авторизация одобрена и~реестр создал
\texttt{merchant\_pending}, а~подтверждение не~пришло (продавец отменил заказ,
не~получил вебхук или забыл подтвердить), правила платёжной системы задают
предельный срок от~авторизации до~клиринга для категории операции; по~его
истечении эквайер не~вправе подтверждать транзакцию, а~эмитент обычно снимает
соответствующий холд
держателя~\cite{visa-core-rules,mc-rules,nspk-mir-rules}.\footnote{По редакции
  Visa Authorization Framework от~13.04.2024 объединённый
  \textit{authorization-to-clearing time frame} составляет 5~дней для всех
  \textit{card-present}-операций и~merchant-initiated transactions, 10~дней для
  cardholder-initiated CNP-операций и~30~дней для \textit{lodging},
  \textit{cruise lines} и~\textit{vehicle rental} при~наличии Estimated
  Authorization Indicator. Mastercard и~НСПК задают собственные пределы; срок
  удержания холда у~эмитента формально его внутренняя политика, на~практике
повторяющая таймфрейм сети.} В~реестре запись \texttt{merchant\_pending}
остаётся неограниченно долго без механизма экспирации. Решением служит фоновый
процесс, сканирующий проводки в~ожидании (\textit{pending}) старше порога
(максимальный срок холда по~категории плюс буфер) и~автоматически порождающий
событие \texttt{auth\_expired}: дебет \texttt{merchant\_pending} и~кредит
\texttt{cardholder\_receivable} обнуляют обязательство. Без этого баланс
\texttt{pending} растёт, искажая отчётность и~блокируя сверку.

\textbf{Гонка возврата и~чарджбэка.} Когда продавец инициирует возврат
на~5\,000~руб. в~тот же день, когда эмитент открывает чарджбэк на~ту же сумму,
оба события дебетуют \texttt{merchant\_available}; без защиты у~продавца
спишется 10\,000~руб. вместо 5\,000. Защитой служит блокировка на~уровне
транзакции: перед записью возврата реестр проверяет, нет ли открытого спора,
а~перед записью чарджбэка~--- нет ли уже проведённого возврата. Если возврат
уже исполнен, чарджбэк переходит в~статус \texttt{already\_refunded},
а~платформа отправляет повторное представление с~доказательством
возврата~\cite{mc-chargeback-guide}; если же чарджбэк пришёл первым, возврат
блокируется, и~оператор разрешает конфликт вручную. Итоговый дебет с~ТСП
в~обоих случаях не~превышает исходную сумму транзакции.

\practicenote{%
  Команде бэк-офиса нужна таблица соответствия внешних идентификаторов
  и~внутренних счетов: авторизация, процессинг, заказ,
  выплата, спор. Без неё даже хороший реестр трудно сверить с~реальным миром.
}

Реестр обязательств опирается на~три инварианта. Двойная запись
из~раздела~\ref{sec:ledger-double-entry} удерживает равенство дебетов
и~кредитов; соответствие счетов реестра и~809-П
из~раздела~\ref{sec:ledger-vs-accounting} сводит продуктовый учёт
с~бухгалтерским по~Плану счетов; идемпотентный приём событий
из~раздела~\ref{sec:ledger-idempotency} удерживает однократность применения.
Когда хотя~бы один из~трёх инвариантов подсчитывается вручную или задним
числом, в~журнале появляются проводки без известного источника, и~расхождения
с~банком разбираются по~переписке бэк-офиса. Реестр обязательств закладывается
с~первой версии: встраивание задним числом обходится в~миграцию данных,
переписывание учётной политики и~пересборку бэк-офиса. Чарджбэки и~orphan auth
из~этой главы~--- проводки, которых нет в~журнале до~момента явной записи;
до~неё они становятся ручной операцией бэк-офиса.
